\documentclass[aps,prd,onecolumn
,tightenlines,letterpaper,notitlepage,
%superscriptaddress,
nofootinbib]{revtex4-1}
\usepackage{amssymb,latexsym}
\usepackage{amsmath,amsbsy,bbm}
\usepackage{epsfig,bm,color}
\usepackage{graphicx,comment}
\unitlength=1mm
\DeclareMathOperator{\st}{str}
\DeclareMathOperator{\Erfc}{Erfc}
\DeclareMathOperator{\Erf}{Erf}
\DeclareMathOperator{\Pf}{Pf}
\DeclareMathOperator{\sign}{sign}

\usepackage{dutchcal}


\usepackage{calligra}

\DeclareMathAlphabet{\mathcalligra}{T1}{calligra}{m}{n}
\DeclareFontShape{T1}{calligra}{m}{n}{<->s*[2.2]callig15}{}
\newcommand{\scriptr}{\mathcalligra{r}\,}
\newcommand{\boldscriptr}{\pmb{\mathcalligra{r}}\,}
% USAGE $\scriptr$ or $\boldscriptr$



%%%%%%%%%%%%%%%%%%%%%%%%%%%%%%%%
\begin{document}
%%%%%%%%%%%%%%%%%%%%%%%%%%%%%%%%


\def\a{{\alpha}}
\def\b{{\beta}}
\def\d{{\delta}}
\def\D{{\Delta}}
\def\X{{\Xi}}
\def\e{{\varepsilon}}
\def\g{{\gamma}}
\def\G{{\Gamma}}
\def\k{{\kappa}}
\def\l{{\lambda}}
\def\L{{\Lambda}}
\def\m{{\mu}}
\def\n{{\nu}}
\def\o{{\omega}}
\def\O{{\Omega}}
\def\S{{\Sigma}}
\def\s{{\sigma}}
\def\th{{\theta}}

\def\ol#1{{\overline{#1}}}

\def\Dslash{D\hskip-0.65em /}
\def\Dtslash{\tilde{D} \hskip-0.65em /}

\def\order{{\mathcal O}}

\def\Dt{{\tilde{D}}}
\def\St{{\tilde{\Sigma}}}

\def\eqref#1{{(\ref{#1})}}

\newcommand{\caf}{\text{\cjRL{b}}}
\newcommand{\he}{${}^4$He}
\newcommand{\hes}{${}^3$He}
\newcommand{\tr}{${}^3$H}
\newcommand{\ls}{\ve{L}\cdot\ve{S}}
\newcommand{\eps}{\epsilon}
\newcommand{\as}{a_s}
\newcommand{\at}{a_t}
\newcommand{\ecm}{E_\textrm{\small c.m.}}
\newcommand{\dq}{\mbox{d\hspace{-.55em}$^-$}}
\newcommand{\mpis}{$m_\pi=137~${\small MeV}}
\newcommand{\mpim}{$m_\pi=450~${\small MeV}}
\newcommand{\mpil}{$m_\pi=806~${\small MeV}}
\newcommand{\muh}{\mu_{^3\text{\scriptsize He}}}
\newcommand{\mut}{\mu_{^3\text{\scriptsize H}}}
\newcommand{\mud}{\mu_\text{\scriptsize D}}
\newcommand{\pode}{\beta_{\text{\scriptsize D},\pm1}}
\newcommand{\poh}{\beta_{^3\text{\scriptsize He}}}
\newcommand{\pot}{\beta_{^3\text{\scriptsize H}}}
\newcommand{\com}[1]{{\scriptsize \sffamily \bfseries \color{red}{#1}}}
\newcommand{\eg}{\textit{e.g.}\;}
\newcommand{\ie}{\textit{i.e.}\;}
\newcommand{\cf}{\textit{c.f.}\;}
\newcommand{\be}{\begin{equation}}
\newcommand{\ee}{\end{equation}}
\newcommand{\la}{\label}
\newcommand{\ber}{\begin{eqnarray}}
\newcommand{\eer}{\end{eqnarray}}
\newcommand{\nn}{\nonumber}
\newcommand{\half}{\frac{1}{2}}
\newcommand{\thalf}{\nicefrac[]{3}{2}}
\newcommand{\bs}[1]{\ensuremath{\boldsymbol{#1}}}
\newcommand{\bea}{\begin{eqnarray}}
\newcommand{\eea}{\end{eqnarray}}
\newcommand{\beq}{\begin{align}}
\newcommand{\eeq}{\end{align}}
\newcommand{\bk}{\bs k}
\newcommand{\bt}{B_{^{3}\text{H}}}
\newcommand{\bh}{B_{^{3}\text{He}}}
\newcommand{\bd}{B_\text{D}}
\newcommand{\ba}{B_\alpha}
\newcommand{\rgm}{$\mathbb{R}$GM}
\newcommand{\ev}[1] {|\bra #1  \ket |^2}
\newcommand{\lam}[1]{$\Lambda=#1~$fm$^{-1}$}
\newcommand{\parg}[1] {\paragraph*{-\,\textit{#1}\,-}}
\newcommand{\nopi}{\pi\hspace{-6pt}/}
\newcommand{\ve}[1]{\ensuremath{\boldsymbol{#1}}}
\newcommand{\xvec}{\bs{x}}
\newcommand{\rvec}{\bs{r}}
\newcommand{\sgve}{\ensuremath{\boldsymbol{\sigma}}}
\newcommand{\tave}{\ensuremath{\boldsymbol{\tau}}}
\newcommand{\na}{\nabla}
\newcommand{\bra}{\langle}
\newcommand{\ket}{\rangle}
\newcommand{\tx}{\tilde{x}}
\newcommand{\eftnopi}{\mbox{EFT($\slashed{\pi}$)}}

\newcommand{\cmment}[1]{\paragraph*{Ecce:\texttt{\textcolor{blue}{#1}}}}

\author{$\mathcal{S}$.~$\mathcal{Mondal}$}
\author{$\mathcal{M}$.~$\mathcal{Elyahu}$}
\author{$\mathcal{N}$.~$\mathcal{Barnea}$}
\author{$\mathcal{J}$.~$\mathcal{Kirscher}$}
 
 
%%%%%%%%%%%%%%%%%%%%%%%%%%%%%%%%%%%%
%%%%%%%%%%%%%%%%%%%%%%%%%%%%%%%%%%%% 
\title{
Strong magnetic fields and contact-interacting few-fermion systems
} 
%%%%%%%%%%%%%%%%%%%%%%%%%%%%%%%%%%%%
%%%%%%%%%%%%%%%%%%%%%%%%%%%%%%%%%%%%


\begin{abstract}
My aim is to predict the critical number of particles which form a self-bound state (droplet)
while their pair interaction is not attractive enough to sustain a bound dimer. As an intermediate
step, I analyze the critical background-field strength for the in-vacuum unstable $A$-body system
to become bound. This field strength shall be analyzed for Bose and Fermi particles.
\end{abstract}

\maketitle

\section*{Scratch Pad}

\begin{align}
\ve{A}&=\frac{1}{2}\,\ve{B}\times\ve{r}\\
&=\frac{1}{2}\,B\rho\,\hat{\phi}\tag{for $\ve{B}=B\hat{z}$}
\intertext{\centering with $\rho=\sqrt{x^2+y^2}$}
\end{align}

\section{Introduction -- what fascinates us?}

Assume a collection of $A$ non-interacting, non-relativistic identical bosons and imagine you could
tune-on a spherically symmetric, attractive pair interaction smoothly. The regime of large enough interaction strength $\nu$
to form a dimer with zero energy and the associated universal 
3-body (original:~\cite{Efimov:1971zz,Efimov:1973awb}; alternative, potentially generalizable explanation:~\cite{AMADO197125,Adhikari:1972par}) and 
cluster spectra (empirical/numerical:~\cite{von_Stecher_2010,PhysRevLett.107.200402,2014FBS....55.1245H,PhysRevA.90.012502,PhysRevA.92.033626,article};
field-theoretical and classical approach:~\cite{Son:2021kkx,Wang:2024qcg})



\section{Theoretical description of strongly interacting fermions in a magnetic
background field}

In order to describe the effect of a static magnetic field, which is not necessarily
perturbatively small, on the low-energy spectra of small nuclei, the 
effective-field-theory Hamiltonian of the latter is augmented by, first, minimally coupling
the derivatives, and second, by the magnetic-moment 1-body interaction (a remnant of the
non-relativistic reduction of the Dirac equation) and the coupling of the $\ve{B}$ field
to the four-nucleon leading-order momentum-independent vertex.

\be\la{eq.hamiltonian}
\hat{H}=
\ee

The so-called Landau spectrum of the minimally coupled, non-interacting limit of this 
Hamiltonian can be derived analytically

\be\la{eq.landau_spect}
E=
\ee

\subsection{Notable features}

\begin{enumerate}
\item Total Spin in not conserved, which means that if we consider a 2-nucleon
system, it resides in a superposition of spin-1 and spin-0 states.
\item Bound states might not have negative eigenvalues because the non-interacting
theory still contains the effect of the magnetic field on the charged particles
and shifts their ground state by, at least, the energy in the lowest Landau level.
\item
\end{enumerate}

\section{Two-nucleon systems}

First, we analyse the predictions in Ref.~\cite{PhysRevLett.116.112301}, \ie, what is
the dependence of the lowest eigenvalues of the 2-proton system on the strength of the
magnetic field?

\section{Three Nucleons}

How does the critical field strength at which the 3-proton system becomes bound compare
with the corresponding strength which is necessary to bind two protons?

\section{Four Nucleons}

Same critical-strength ratio? Now we extrapolate and predict that no more than 
$\vert\ve{B}\vert\sim 4~$T is needed to bind an $A$-proton cluster.

\bibliography{contacts-in-strong-B}

\end{document}